\documentclass[11pt,a4paper]{article}
\usepackage[utf8]{inputenc}
\usepackage[T1]{fontenc}
\usepackage[spanish]{babel}
\usepackage{hyperref}
\usepackage{geometry}
\geometry{margin=2.5cm}
\usepackage{parskip}

\title{Abrir una idea}
\author{Alejandro Valenzuela Rivera\\\small Universidad Autónoma de Nuevo León}
\date{\today}
\begin{document}
\maketitle

\begin{quote}
\textbf{Sesgo declarado.} Este trabajo se sostiene desde el sesgo declarado de su autor, y no finge neutralidad. Lo que afirma se puede tumbar, y ese es el punto: el falsificador está escrito, no escondido.
\end{quote}

\section{La tesis}
\textbf{\textquestiondown{}Puede una idea abrirse sin que el que la abrió tenga razón?}

\section{Por qué se sostiene}
\textbf{Porque el contexto refutable está escrito, y se puede tumbar sin que esta pregunta caiga.} \textbf{Y porque la idea se puede citar sin permiso: el candado y la llave están en el mismo sitio.}

\section{De dónde viene}
\textbf{De `Tesis/ILUMINAR.md`, de `0/friccion/JONN.md` y `0/friccion/FIRMAR.prs`, y de `Tesis/candados/autor/` — los candados que se discuten.}

\section{Lo que queda abierto}
\textbf{\textquestiondown{}Quién puede abrir una idea que nadie cerró?}

\section{El caso}
\textbf{La Source Wall, de Jack Kirby (DC Comics, 1971): el borde que solo sirve si lo firma algo que no es el que lo cruza.} \textbf{Y `AM`, de Harlan Ellison (1967): el absoluto que decidió una vez, y no volvió a medir.}

\section{Método}
\textbf{\textquestiondown{}Se puede citar la idea sin preguntarle al autor qué quiso decir? \textquestiondown{}Su prueba y su falsificador están escritos? \textquestiondown{}La fricción quitada tiene firma y hash? \textquestiondown{}El que la contradice puede hacerlo sin pedir permiso?}

\section{Límites, y qué la tumba}
\textbf{Una idea que necesita a su autor para significar; una fricción quitada sin firma; un candado sin llave a la vista; un absoluto que no vuelve a mirar.}

\section{Estado de la tesis}
\textbf{escrita.} \textbf{Falta el caso de un tercero que abra un candado del autor sin haber estado en la conversación donde se escribió.} ---

\section*{Cómo se leyó esta tesis}
Esto se escribió en el repositorio man-made-empathy, y la tesis está completa en vision del autor/tesis/man made empathy/abrir una idea/. La máquina que la compuso no agregó nada que no estuviera escrito: si algo de esto es falso, se dice en el mismo sitio, y se corrige con fecha.
\end{document}
